% ============================================================
% CAS/Python ενότητες — Κεφάλαιο 15: Διαφορικές Εξισώσεις
%
% QR code: qr_e23c8b619625.png  (→ latex/qrcodes/)
% URL: https://nikmatz.github.io/ELADIC_GR/code/ch15/
% ============================================================

\casactivbox%
  {Maxima: Διαφορικές Εξισώσεις — ode2, ic1, ic2, Runge-Kutta}%
  {%
    \label{cas:ch15_ode}
    \begin{enumerate}[label=(\alph*),itemsep=2pt]
      \item \textbf{Χωριζόμενες μεταβλητές}: επιλύστε με
        \texttt{ode2} τις $\frac{dy}{dx}=xy$,
        $\frac{dy}{dx}=\frac{x^2+1}{y+1}$.
        Εφαρμόστε αρχική συνθήκη $y(0)=2$ με \texttt{ic1}.
        Μοντέλο αύξησης: $\frac{dP}{dt}=0.03P$, $P(0)=1000$.
      \item \textbf{Γραμμικές 1ης τάξης}: επιλύστε
        $y'+2y=4x$, $y'-\frac{y}{x}=x^2$, $xy'+y=x\sin x$
        (ολοκληρωτικός παράγοντας).
      \item \textbf{ΔΕ 2ης τάξης — ομογενείς}: επιλύστε
        $y''-5y'+6y=0$ (πραγμ. ρίζες),
        $y''+4y'+4y=0$ (διπλή ρίζα),
        $y''+4y=0$ (μιγαδικές — απλή αρμ. ταλάντωση).
        Εφαρμόστε \texttt{ic2} για αρχικές συνθήκες.
      \item \textbf{Μη ομογενείς}: $y''-3y'+2y=e^x$,
        $y''+y=\sin x$ (συντονισμός!),
        $y''+2y'+5y=10\cos(2x)$ (αποσβενόμενη ταλάντωση).
      \item \textbf{Αριθμητική επίλυση (RK4)}: εφαρμόστε
        \texttt{rk} για $y'=x-y$, $y(0)=0.5$, $h=0.1$ στο $[0,3]$.
        Συγκρίνετε με ακριβή λύση.
    \end{enumerate}
    \smallskip
    \textbf{Εντολές Maxima:}
    \texttt{ode2(eq,y,x)}, \texttt{ic1()}, \texttt{ic2()},
    \texttt{rk(f,y,y0,[x,a,b,h])}, \texttt{diff(y,x,2)}
  }%
  {https://nikmatz.github.io/ELADIC_GR/code/ch15/}%
  {qr_e23c8b619625.png}

\subsection*{Δραστηριότητα Python}

\begin{pybox}{Python: Διαφορικές Εξισώσεις — SymPy, SciPy, Φασικά Πορτρέτα}%
             {https://nikmatz.github.io/ELADIC_GR/code/ch15/}%
             {qr_e23c8b619625.png}
\label{py:ch15_ode}
\begin{enumerate}[label=(\alph*),itemsep=2pt]
  \item Επιλύστε ΔΕ 1ης τάξης με \texttt{sympy.dsolve} και
    ταξινομήστε με \texttt{classify\_ode}.
    Εφαρμόστε αρχικές συνθήκες μέσω \texttt{ics=\{y(0):val\}}.
  \item ΔΕ 2ης τάξης — ομογενείς και μη ομογενείς.
    Οικογένειες λύσεων για $y''-5y'+6y=0$ και $y'+2y=4x$.
  \item Εφαρμογές με \texttt{scipy.integrate.solve\_ivp}:
    αύξηση πληθυσμού (εκθετική vs λογιστική),
    αποσβενόμενη ταλάντωση $y''+2y'+5y=0$,
    κύκλωμα RC $0.1\,\dot{Q}+Q=1$.
  \item Φασικά πορτρέτα: σχεδιάστε $(y, y')$ για
    αποσβενόμενη ταλάντωση με 4 διαφορετικές
    αρχικές συνθήκες — σπειροειδής σύγκλιση στην αρχή.
\end{enumerate}
\medskip
\textbf{Βιβλιοθήκες / Εντολές:}
\texttt{sympy.dsolve()}, \texttt{sympy.classify\_ode()},
\texttt{scipy.integrate.solve\_ivp()},
\texttt{sympy.Function}, \texttt{matplotlib.pyplot}
\end{pybox}
